\documentclass[twocolumn]{aastex631}
\begin{document}
\title{An age dependence of the Galactic alpha-knee across galactocentric radius}
\author{NebulaMind Lab (autonomous pipeline)}
\affiliation{NebulaMind Lab, \url{https://lab.nebulamind.net}}
\begin{abstract}
Age-resolved alpha-knee from an APOGEE (DR18) 3-table join of 329,785 giants (apogeeDistMass C/N ages + distances, apogeeStar Galactic coordinates, aspcapStar [Fe/H]-[O/Fe]). The [Fe/H]\_knee is located per (R\_g x age) cell with a broken-line ridge fit (bootstrap error) in 32 populated cells; median [Fe/H]\_knee = -0.57. Gradients: d[Fe/H]\_knee/d(age)|\_R\_g = +0.0140+/-0.0086 dex/Gyr; d[Fe/H]\_knee/dR\_g|\_age = +0.0215+/-0.0120 dex/kpc. Non-circularity: the age-gradient sign HOLDS under the abundance-scale swap ([Fe/H]->[M/H], [O/Fe]->[alpha/M]). Generated autonomously from public data (apogee) via the NebulaMind Lab runner: a bounded, reproducible, descriptive study.
\end{abstract}
\section{Introduction}
The study of the Milky Way's chemical evolution has been a longstanding area of interest in astrophysics. Previous research has explored various aspects of this topic, including the use of rotation-based ages for APOGEE-Kepler cool dwarf stars [Claytor2020] and the identification of young alpha-rich stars in different Galactic regions [Grisoni2024]. Additionally, studies have focused on stellar model tests and age determination for RGB stars from the APO-K2 catalogue [Valle2024] and intermediate-age alpha-rich populations in K2 [Warfield2021]. Building upon these works, we aim to investigate the age-resolved alpha-knee using data from the APOGEE DR18 catalog.
\section{Data and method}
To achieve this goal, we employ a methodology that involves pulling three tables (apogeeDistMass, apogeeStar, aspcapStar) from the SDSS DR18 APOGEE catalog via SkyServer raw-HTTP. We chunk the data by sky region and apply pacing, retry/backoff, and disk cache to manage the process efficiently. The tables are joined in-process on APOGEE\_ID, and flag/quality cuts (STAR\_BAD, per-element flags, SNR, abundance errors, 1-14 Gyr ages) are applied using Python. We compute R\_g from glon/glat/distance with R0=8.122 kpc and utilize spectroscopic C/N ages while acknowledging the known systematics associated with them.
\section{Result}
Our analysis yields an age-resolved alpha-knee measurement based on a sample of 329,785 giants. The [Fe/H]\_knee is determined per (R\_g x age) cell using a broken-line ridge fit with bootstrap error in 32 populated cells, resulting in a median [Fe/H]\_knee value of -0.57. We find gradients of d[Fe/H]\_knee/d(age)|\_R\_g = +0.0140+/-0.0086 dex/Gyr and d[Fe/H]\_knee/dR\_g|\_age = +0.0215+/-0.0120 dex/kpc. Notably, the age-gradient sign remains consistent under a non-circularity test that swaps the abundance calibration scale.
\begin{figure}\centering\includegraphics[width=\columnwidth]{result.png}\caption{An age dependence of the Galactic alpha-knee across galactocentric radius}\end{figure}
\section{Caveats}
It is essential to acknowledge the limitations of our approach. The automated nature of this study may introduce biases or overlook subtle nuances in the data. Additionally, relying on a single selection criterion and uncalibrated measurements can lead to potential inaccuracies. The spectroscopic C/N ages used in this analysis are subject to known systematics, which may impact the reliability of our results. Furthermore, the absence of external validation or comparison with other datasets restricts the generalizability of our findings. These caveats highlight the need for further refinement and verification through more comprehensive and diverse analyses.
\section*{References}
\footnotesize
[Claytor2020] Claytor et al. (2020). Chemical Evolution in the Milky Way: Rotation-based Ages for APOGEE-Kepler Cool Dwarf Stars. bibcode 2020ApJ...888...43C \par [Grisoni2024] Grisoni et al. (2024). K2 results for "young" α-rich stars in the Galaxy. bibcode 2024A\&A...683A.111G \par [Valle2024] Valle et al. (2024). Stellar model tests and age determination for RGB stars from the APO-K2 catalogue. bibcode 2024A\&A...690A.323V \par [Warfield2021] Warfield et al. (2021). An Intermediate-age Alpha-rich Galactic Population in K2. bibcode 2021AJ....161..100W

\end{document}
